% =====================================================================
%  CHAPITRE 16 — COMPRÉHENSION ORALE
% =====================================================================
\chapter{Comprendre l'anglais parlé : la méthode d'écoute}

\begin{quote}\itshape
« Je comprends quand je lis, mais dès qu'on me parle, tout disparaît. »
C'est la plainte la plus fréquente --- et elle a une cause précise, donc une solution précise.
\end{quote}

\section{Pourquoi vous ne comprenez pas (les cinq causes)}
\begin{enumerate}
  \item \textbf{Vous cherchez des mots isolés.} Le natif ne dit pas les mots séparément :
  \en{What do you want to do?} devient \pron{oua-da-ya-wou-na-dou}. Votre oreille cherche
  cinq mots français-style ; il n'y en a qu'un seul, fondu. (Voir le chapitre 6.)
  \item \textbf{Votre oreille n'est pas entraînée aux sons manquants.} \ipatheta, \ipadh, /h/,
  /w/, \ipang, \ipeschwa : tant que vous ne les distinguez pas, des dizaines de mots
  courants vous échappent.
  \item \textbf{Le niveau est trop élevé.} Écouter CNN au mois 2 ne fait pas progresser :
  cela démoralise. Il faut 90\,\% de compréhension pour apprendre.
  \item \textbf{Vous écoutez passivement.} Le fond sonore aide au rythme, mais ne crée pas
  de compétence. Seule l'écoute \textbf{active} (avec tâche) progresse.
  \item \textbf{Vous traduisez en français.} Traduire prend 2 secondes ; la phrase suivante
  arrive avant. Il faut comprendre \textbf{directement} --- c'est un entraînement, pas un don.
\end{enumerate}

\section{Les deux écoutes : intensive et extensive}
\begin{tableau}[Les deux régimes, tous deux indispensables]
\begin{tabularx}{\linewidth}{@{}p{3.2cm}p{5.0cm}X@{}}
\entete{} & \entete{Écoute intensive} & \entete{Écoute extensive} \tabularnewline
\hline
Durée & 15 à 25 min/jour & 30 à 60 min/jour (temps mort) \tabularnewline
Matériau & 1 à 3 minutes d'audio, niveau 90\,\% compris & Contenu long, facile, agréable \tabularnewline
Méthode & 5 passes, dictée, shadowing, analyse & On écoute en faisant autre chose, sans pause \tabularnewline
Objectif & Précision : sons, liaisons, mots nouveaux & Volume : rythme, familiarité, endurance \tabularnewline
Effet & Progresse vite, fatigue & Progresse lentement, ne fatigue pas \tabularnewline
\end{tabularx}
\end{tableau}

\subsection{Le protocole des 5 passes (écoute intensive, 20 minutes)}
\begin{enumerate}
  \item \textbf{Passe 1 --- écoute aveugle (2 min).} Écoutez l'audio sans texte, deux fois.
  Notez seulement : quel est le sujet ? qui parle ? quel est le ton ? Ne cherchez pas à tout
  comprendre.
  \item \textbf{Passe 2 --- écoute avec le texte (3 min).} Lisez la transcription en écoutant.
  Soulignez les mots inconnus et les endroits où le son ne correspondait pas à ce que vous
  attendiez.
  \item \textbf{Passe 3 --- écoute sans le texte (3 min).} Deux fois. Vous devez maintenant
  comprendre 95\,\%. Sinon, revenez à la passe 2.
  \item \textbf{Passe 4 --- dictée (7 min).} Écrivez une phrase sur cinq, phrase par phrase,
  avec pauses. Comparez avec le texte. \textbf{Analysez chaque erreur} : mot inconnu ?
  liaison non perçue ? son non reconnu ? C'est l'étape qui fait le plus progresser.
  \item \textbf{Passe 5 --- shadowing (5 min).} Répétez en même temps que l'audio, sans le
  texte, en imitant tout.
\end{enumerate}

\begin{astuce}[Le même audio trois jours de suite]
Jour 1 : protocole complet. Jour 2 : passe 3 + passe 5 seulement (7 min). Jour 3 : écoute
simple en marchant. Ce triple contact ancre l'audio dans votre mémoire --- vous retiendrez
80\,\% des expressions, contre 20\,\% avec un audio écouté une fois.
\end{astuce}

\section{La méthode « série » : 3 passes par épisode}
\begin{tableau}[Comment travailler une série sans se décourager]
\begin{tabularx}{\linewidth}{@{}p{2.4cm}p{3.6cm}X@{}}
\entete{Passe} & \entete{Sous-titres} & \entete{Objectif} \tabularnewline
\hline
1 & Français (mois 1-3) puis anglais (mois 4+) & Comprendre l'histoire, prendre du plaisir. Sans plaisir, on abandonne. \tabularnewline
2 & \textbf{Anglais} & Relier le son à l'écrit. Vous découvrirez que vous « saviez » la moitié des mots sans les reconnaître à l'oreille. \tabularnewline
3 & Aucun & Tester votre oreille réelle. 10 à 15 minutes suffisent --- le reste de l'épisode en passe 1 ou 2. \tabularnewline
\end{tabularx}
\end{tableau}

\begin{memo}[Le choix des séries]
Préférez les \textbf{sitcoms du quotidien} (dialogues courts, situations répétitives, humour
simple) aux séries policières ou fantastiques (vocabulaire technique, débit rapide, bruit de
fond). Regardez \textbf{deux fois le même épisode} : la seconde fois, tout devient audible,
et c'est cette sensation qui vous prouve vos progrès.
\end{memo}

\section{La méthode « chanson » (15 minutes, une fois par semaine)}
\begin{enumerate}
  \item Choisissez une chanson au débit clair (pop, folk, soul --- évitez le rap au début).
  \item Écoutez deux fois sans paroles, en notant les mots que vous reconnaissez.
  \item Lisez les paroles en écoutant : soulignez le vocabulaire nouveau et les liaisons.
  \item Cherchez le sens global (de quoi parle la chanson ? quelle émotion ?).
  \item Chantez --- oui, vraiment. Chanter installe le rythme, la longueur des voyelles et
  les enchaînements mieux que n'importe quel exercice.
  \item Ajoutez 5 expressions de la chanson à votre Anki.
\end{enumerate}

\section{Dix textes d'écoute à travailler}
Les textes ci-dessous sont conçus pour être \textbf{écoutés} : faites-les lire par une
application de synthèse vocale (votre téléphone, un dictionnaire en ligne, un traducteur
vocal), par un partenaire, ou lisez-les à voix haute en vous enregistrant, puis travaillez-les
sans le texte. Chaque texte est suivi de questions ; les réponses sont en fin de section.

\subsection{Texte 1 --- niveau A1 : se présenter}
\begin{dialoguebox}[Texte 1 --- à écouter 5 fois]
\begin{anglais}
Hello! My name is Awa Diop. I'm twenty-four years old and I come from Senegal. I live in
Thiès with my parents and my two brothers. I work in a small shop near the market. I open
the shop at eight in the morning and I close it at six in the evening. In my free time, I
like listening to music and playing basketball with my friends. I'm learning English because
I want to work with customers from other countries. My favourite subject is English, but I
find pronunciation difficult. On Sundays, I usually visit my grandmother and we eat rice
with fish together.
\end{anglais}
\end{dialoguebox}
\textbf{Questions :} 1. Quel âge a Awa ? 2. Où habite-t-elle et avec qui ?
3. Quelles sont ses heures de travail ? 4. Pourquoi apprend-elle l'anglais ?
5. Que fait-elle le dimanche ?

\subsection{Texte 2 --- niveau A2 : une journée de travail}
\begin{dialoguebox}[Texte 2 --- à écouter 5 fois]
\begin{anglais}
I usually get up at six, but yesterday was different. I woke up at nine because I had worked
very late the night before. When I looked at my phone, I saw five missed calls from my boss.
I called him back immediately and he wasn't angry --- he just needed the report for the
morning meeting. I took a taxi because I didn't have time to wait for the bus, and the
traffic was terrible, so I arrived at eleven. In the end, I finished the report at two
o'clock and my boss thanked me. That evening, I went to bed at nine. I've never been so
tired in my life.
\end{anglais}
\end{dialoguebox}
\textbf{Questions :} 1. Pourquoi s'est-il réveillé à neuf heures ? 2. Qu'a-t-il vu sur son
téléphone ? 3. Son patron était-il fâché ? 4. Pourquoi a-t-il pris un taxi ?
5. À quelle heure a-t-il fini le rapport ? 6. Repérez les quatre temps utilisés.

\subsection{Texte 3 --- niveau A2 : au restaurant}
\begin{dialoguebox}[Texte 3 --- à écouter 5 fois]
\begin{anglais}
Good evening! Do you have a reservation? No? That's not a problem --- we have a table by the
window. Here are the menus. Today's special is grilled fish with vegetables, and we also
have chicken with rice. Would you like something to drink while you decide? We have fresh
orange juice, mineral water, and mint tea. I'll give you a few minutes. By the way, the fish
is very fresh --- it came from the market this morning. Let me know when you're ready to
order, and if you have any allergies, please tell me now.
\end{anglais}
\end{dialoguebox}
\textbf{Questions :} 1. Le client avait-il réservé ? 2. Où est placée la table ?
3. Quels sont les deux plats proposés ? 4. Quelles boissons sont citées ?
5. Pourquoi le poisson est-il « very fresh » ?

\subsection{Texte 4 --- niveau B1 : un appel téléphonique professionnel}
\begin{dialoguebox}[Texte 4 --- à écouter 5 fois]
\begin{anglais}
Good morning, West Africa Logistics, Fatou speaking. How can I help you? \ldots\ I see. So
you're calling about the container that was supposed to arrive on Tuesday? Let me check that
for you. Could you give me the reference number, please? Thank you. Ah, yes --- I can see
the problem. Unfortunately, the shipment has been delayed because of bad weather at the port,
and it's now expected on Friday. I do apologise for the inconvenience. What I can do is send
you a confirmation by email within the hour, and I'll also ask my colleague to call you if
there's any further change. Would that be all right? \ldots\ Perfect. Is there anything else
I can help you with today? Thank you for calling, Mr Ba. Have a good day.
\end{anglais}
\end{dialoguebox}
\textbf{Questions :} 1. Pourquoi le client appelle-t-il ? 2. Quelle est la cause du retard ?
3. Quand le conteneur doit-il arriver ? 4. Quelles sont les deux promesses de Fatou ?
5. Relevez trois formules de politesse professionnelle.

\subsection{Texte 5 --- niveau B1 : donner son opinion}
\begin{dialoguebox}[Texte 5 --- à écouter 5 fois]
\begin{anglais}
In my opinion, social media has changed the way young people learn English, and most of the
change is positive. Ten years ago, if you wanted to hear real English, you had to find a
teacher or a cassette. Today, you can watch a film, follow a footballer, or listen to a
podcast on your phone --- for free, and as often as you like. On the other hand, I think
there's a danger. Many young people believe that watching videos is enough, and they never
open their mouths. Understanding is not speaking. So my advice is simple: use social media
every day, but spend at least twenty minutes a day producing English --- out loud, alone if
necessary. That's how you turn input into ability.
\end{anglais}
\end{dialoguebox}
\textbf{Questions :} 1. Quelle est l'opinion générale du locuteur ? 2. Quelle comparaison
fait-il avec il y a dix ans ? 3. Quel danger identifie-t-il ? 4. Que signifie pour lui
« Understanding is not speaking » ? 5. Quel conseil donne-t-il précisément ?

\subsection{Texte 6 --- niveau B1 : raconter un voyage}
\begin{dialoguebox}[Texte 6 --- à écouter 5 fois]
\begin{anglais}
Last year I travelled to Ghana for the first time, and honestly, I wasn't expecting to enjoy
it so much. I'd planned everything: the flight, the hotel, the places I wanted to visit.
But of course, nothing went as planned. On the second day, while I was waiting for a bus to
Cape Coast, I started talking to a man called Kwame. He told me he was going to the same
place, and that his cousin had a car. Two hours later, I was sitting in the front seat,
eating bread with eggs, and listening to stories about his family. We spent three days
together. I've kept in touch with him ever since --- we speak on the phone almost every
month. If I hadn't missed that bus, I'd never have met him. Travelling isn't about the
places; it's about the people you meet in them.
\end{anglais}
\end{dialoguebox}
\textbf{Questions :} 1. Où le narrateur est-il allé ? 2. Qu'est-ce qui ne s'est pas passé
comme prévu ? 3. Comment a-t-il rencontré Kwame ? 4. Sont-ils encore en contact ?
5. Repérez le conditionnel 3 et le \en{present perfect} dans le texte.
6. Quelle est la conclusion du narrateur ?

\subsection{Texte 7 --- niveau B2 : les actualités}
\begin{dialoguebox}[Texte 7 --- à écouter 5 fois]
\begin{anglais}
Good evening. The government announced this afternoon that a new programme will provide
free digital skills training to fifty thousand young people over the next three years.
Speaking at a press conference in the capital, the Minister of Youth said that the scheme
had already been tested successfully in two regions, where more than seventy per cent of
participants found work within six months. Critics, however, have questioned whether the
funding will be sufficient. One economist told our correspondent that similar initiatives
in the past had failed to reach rural areas. The programme is expected to begin in January,
and registration will open online next week. In other news, the national team has qualified
for the Africa Cup of Nations after a two-nil victory last night.
\end{anglais}
\end{dialoguebox}
\textbf{Questions :} 1. Combien de jeunes seront formés, et en combien de temps ?
2. Quel résultat a été observé dans les deux régions test ? 3. Quelle critique est formulée ?
4. Quand le programme commencera-t-il ? 5. Quelle est la seconde nouvelle ?
6. Relevez deux exemples de voix passive.

\subsection{Texte 8 --- niveau B2 : un entretien}
\begin{dialoguebox}[Texte 8 --- à écouter 5 fois]
\begin{anglais}
So, tell me --- what would you say is your greatest professional achievement? \ldots\ Hmm.
I think it would have to be the reorganisation of our accounting system two years ago.
Before that, we were closing the monthly accounts by hand, and it took almost two weeks.
I suggested that we move to a new software, but the idea wasn't popular at first --- people
were used to the old way, and nobody wanted extra work. So I did two things: I trained the
team myself, in the evenings, and I showed the management how much time we would save.
Six months later, we were closing the accounts in three days instead of fourteen. To be
honest, I'm prouder of the way I convinced people than of the result itself. Getting people
to change is harder than changing a system.
\end{anglais}
\end{dialoguebox}
\textbf{Questions :} 1. Quelle question est posée au début ? 2. Quel problème existait avant ?
3. Pourquoi l'idée n'était-elle pas populaire ? 4. Quelles sont les deux choses faites par
le candidat ? 5. Quel est le résultat chiffré ? 6. De quoi est-il le plus fier, et pourquoi ?

\subsection{Texte 9 --- niveau A2/B1 : conversation familière}
\begin{dialoguebox}[Texte 9 --- à écouter 5 fois (débit rapide, réductions)]
\begin{anglais}
Hey, what's up? Long time no see! \ldots\ Yeah, I know, I've been really busy lately. Work's
been crazy --- we've got a new manager and he wants everything done yesterday. Anyway, how
about you? \ldots\ Oh, that's great! Congrats! When do you start? \ldots\ Nice. Hey, we
should catch up properly sometime --- maybe grab a coffee this weekend? I'm free Saturday
afternoon. \ldots\ Cool, let's say three o'clock at the place near the beach? Perfect.
I gotta run now, but I'll text you the details. Take care!
\end{anglais}
\end{dialoguebox}
\textbf{Questions :} 1. Quel est le registre (formel ou familier) ? 2. Que veut dire
\en{everything done yesterday} ? 3. Quelle bonne nouvelle l'autre personne annonce-t-elle ?
4. Que proposent-ils de faire ? 5. Où et quand se retrouveront-ils ?
6. Relevez quatre réductions familières.

\subsection{Texte 10 --- niveau B2/C1 : conférence courte}
\begin{dialoguebox}[Texte 10 --- à écouter 5 fois]
\begin{anglais}
Here's a question I'd like you to sit with for a moment: why do most adults find language
learning so much harder than children do? The usual answer is that children's brains are
more flexible. That's partly true, but I'd like to suggest something else. Children are
allowed to be terrible. A three-year-old says ``I goed to school'' and everybody smiles and
repeats the correct form. Nobody laughs at them; nobody thinks they're stupid. Adults, on
the other hand, have a reputation to protect. We're afraid of sounding foolish, so we stay
silent --- and silence is the one thing that guarantees you'll never improve. So my advice
isn't about grammar or vocabulary. It's much simpler: give yourself permission to be bad at
something, publicly, every single day, for a year. That's the whole secret.
\end{anglais}
\end{dialoguebox}
\textbf{Questions :} 1. Quelle est la question posée ? 2. Quelle est la réponse habituelle,
et pourquoi est-elle insuffisante ? 3. Quel avantage les enfants ont-ils selon la conférencière ?
4. Que craignent les adultes ? 5. Quel est le conseil final ? 6. Expliquez
\en{sit with a moment} et \en{give yourself permission}.

\subsection{Corrigés des dix textes}
\begin{corrige}[Textes 1 à 10 --- réponses essentielles]
\textbf{1.} 24 ans ; Thiès, avec ses parents et ses deux frères ; 8 h - 18 h ;
pour travailler avec des clients étrangers ; elle visite sa grand-mère et ils mangent du
riz au poisson.\\
\textbf{2.} Il avait travaillé très tard la veille ; cinq appels manqués de son patron ;
non, il avait seulement besoin du rapport ; pas le temps d'attendre le bus (et circulation
terrible) ; à 14 h ; passés utilisés : \en{present simple} (\en{I usually get up}),
\en{past simple} (\en{I woke up}), \en{past perfect} (\en{I had worked}),
\en{present perfect} (\en{I've never been}).\\
\textbf{3.} Non ; près de la fenêtre ; poisson grillé avec légumes, poulet avec riz ;
jus d'orange frais, eau minérale, thé à la menthe ; il vient du marché ce matin.\\
\textbf{4.} Un conteneur qui devait arriver mardi ; le mauvais temps au port ; vendredi ;
confirmation par e-mail dans l'heure + appel du collègue en cas de changement ;
\en{I do apologise for the inconvenience}, \en{Would that be all right?},
\en{Is there anything else I can help you with today?}\\
\textbf{5.} Plutôt positive ; avant, il fallait un professeur ou une cassette ;
croire que regarder suffit et ne jamais parler ; comprendre ne remplace pas produire ;
utiliser les réseaux sociaux chaque jour \textbf{et} produire 20 minutes à voix haute.\\
\textbf{6.} Au Ghana ; rien ne s'est passé comme prévu / il a raté le bus ; en attendant le
bus pour Cape Coast ; oui, ils se parlent presque chaque mois ;
\en{If I hadn't missed that bus, I'd never have met him} (conditionnel 3) et
\en{I've kept in touch with him ever since} (\en{present perfect}) ;
le voyage compte par les personnes rencontrées.\\
\textbf{7.} 50\,000 jeunes en trois ans ; plus de 70\,\% ont trouvé un emploi en six mois ;
le financement pourrait être insuffisant et les zones rurales oubliées ; en janvier,
inscriptions en ligne la semaine prochaine ; l'équipe nationale s'est qualifiée
(victoire 2-0) ; passifs : \en{the scheme had already been tested}, \en{the programme is expected to begin}.\\
\textbf{8.} « Quelle est votre plus grande réalisation professionnelle ? » ; la clôture des
comptes prenait deux semaines, à la main ; les gens étaient habitués et ne voulaient pas de
travail supplémentaire ; il a formé l'équipe lui-même le soir + a montré les gains de temps
à la direction ; trois jours au lieu de quatorze ; il est plus fier d'avoir convaincu les
gens, car changer les habitudes est plus dur que changer un système.\\
\textbf{9.} Familier ; tout doit être fait immédiatement (hier = déjà trop tard) ;
l'autre a un nouveau travail (ou une bonne nouvelle personnelle) ; prendre un café samedi ;
à 15 h, près de la plage ; réductions : \en{what's up}, \en{I've been} \pron{aïv-bin}, \en{we've got}, \en{gotta run}, \en{Let's say}, \en{I'll text you}.\\
\textbf{10.} Pourquoi les adultes apprennent-ils moins facilement que les enfants ?
La flexibilité cérébrale --- vraie mais incomplète ; les enfants ont le droit d'être
mauvais, personne ne se moque ; les adultes craignent d'avoir l'air idiots et se taisent ;
se donner la permission d'être mauvais, en public, chaque jour, pendant un an ;
\en{sit with} = garder une question en tête, y réfléchir ; \en{give yourself permission}
= s'autoriser.
\end{corrige}

\section{Le plan d'écoute sur 12 mois}
\begin{tableau}[Ce que vous écoutez, mois par mois]
\begin{tabularx}{\linewidth}{@{}p{1.8cm}p{4.6cm}p{4.2cm}X@{}}
\entete{Mois} & \entete{Intensif (20 min/jour)} & \entete{Extensif (temps mort)} & \entete{Objectif} \tabularnewline
\hline
1-2 & Dialogues de ce livre lus à voix haute + audios pour débutants (vitesse 0,75×) & Chansons lentes, radio en fond & Reconnaître les 44 sons et le rythme \tabularnewline
3-4 & Textes 1 à 5 ci-dessus, audios A2, dessins animés & Podcasts pour apprenants, séries avec sous-titres français & Comprendre un dialogue simple sans texte \tabularnewline
5-6 & Textes 6 à 10, podcasts A2/B1, interviews courtes & Sitcoms avec sous-titres anglais & Comprendre 70\,\% d'une sitcom \tabularnewline
7-8 & Actualités lentes, conférences courtes, TED-Ed & Séries en sous-titres anglais, documentaires & Suivre une conférence simple \tabularnewline
9-10 & Journaux télévisés, podcasts natifs, débats & Films, humour, interviews longues & Comprendre les informations à 80\,\% \tabularnewline
11-12 & Contenu natif non adapté, accents variés & Films sans sous-titres, radio, humour & Comprendre un film sans sous-titres \tabularnewline
\end{tabularx}
\end{tableau}

\begin{mission}[Chapitre 16 --- la semaine d'écoute type]
\textbf{Lundi :} texte 1, protocole des 5 passes. \textbf{Mardi :} texte 1, passes 3 et 5
seulement + texte 2, passes 1-2. \textbf{Mercredi :} texte 2 complet.
\textbf{Jeudi :} dictée de 8 phrases d'un texte déjà travaillé. \textbf{Vendredi :} chanson
(méthode chanson complète). \textbf{Samedi :} un épisode de série en 3 passes.
\textbf{Dimanche :} réécoute passive de tout le matériel de la semaine en marchant.
\end{mission}
